\clearpage
\begin{examplebox}
\textbf{例1（基础母题：给方程，证明相交处正交）。}

已知
\[
E:\ \frac{x^2}{16}+\frac{y^2}{12}=1,\qquad
H:\ x^2-\frac{y^2}{3}=1.
\]
求全部公共点，并证明每个交点处两曲线的切线相互垂直。

\textbf{识别信号：}本例 $a^2=16,b^2=12,A^2=1,B^2=3$，有
$a^2-A^2=16-1=15=b^2+B^2=12+3$，\textbf{一眼确认共焦}；
若题目要求展示计算，再求斜率。

\textbf{解：}联立 $x^2=4,y^2=9$，四个交点为
\[
\boxed{(2,3),(2,-3),(-2,3),(-2,-3)}.
\]
两条切线的斜率为 $k_E=-3x/(4y)$、$k_H=3x/y$，
所以在四个交点都满足
\[
k_Ek_H=-\frac{9x^2}{4y^2}
=-\frac{9\cdot4}{4\cdot9}=-1.
\]
\textbf{结果：}四个交点处全部正交。

\textbf{直观数值核验：}在 $P=(2,3)$，
$t_E:x+2y=8$，$t_H:2x-y=1$，斜率为 $-1/2$ 和 $2$。
\end{examplebox}

\clearpage
\begin{examplebox}
\textbf{例2（迁移母题：只给焦距与轴长，不给公共点）。}

已知共焦的椭圆半长轴为 $a=5$，双曲线实半轴为 $A=2$，
公共焦半距 $c=3$。求所有交点坐标，并写出第一象限交点处的切线方程。

\textbf{解：}由 $c^2=a^2-b^2=A^2+B^2=9$，得
$b^2=25-9=16$、$B^2=9-4=5$，并核对
$a^2-A^2=25-4=21=b^2+B^2=16+5$。故椭圆、双曲线分别为
\[
\frac{x^2}{25}+\frac{y^2}{16}=1,\qquad
\frac{x^2}{4}-\frac{y^2}{5}=1.
\]
由统一的交点坐标式 $x^2=a^2A^2/c^2$、$y^2=b^2B^2/c^2$，
\[
x^2=\frac{25\cdot4}{9}=\frac{100}{9},\qquad
 y^2=\frac{16\cdot5}{9}=\frac{80}{9}.
\]
四交点为 $\left(\pm\frac{10}{3},\pm\frac{4\sqrt5}{3}\right)$。
第一象限点 $P=(10/3,4\sqrt5/3)$，两切线为
\[
\frac{2x}{15}+\frac{\sqrt5\,y}{12}=1,\qquad
\frac{5x}{6}-\frac{4\sqrt5\,y}{15}=1.
\]
对应法向量点积为
\[
\frac{2}{15}\cdot\frac56
+\frac{\sqrt5}{12}\cdot\left(-\frac{4\sqrt5}{15}\right)
=\frac19-\frac19=0.
\]
\textbf{结果：}这两条切线垂直，而整个求解过程不需要计算焦角。
\end{examplebox}

\clearpage
\begin{examplebox}
\textbf{例3（辨析：焦角不等于切线夹角）。}

仍取例1的曲线，设 $P=(2,3)$，求 $\angle F_1PF_2$，
并与椭圆、双曲线切线的夹角比较。

\textbf{解：}焦点 $F_1=(-2,0),F_2=(2,0)$，
所以 $PF_1=5,PF_2=3,F_1F_2=4$，由余弦定理
\[
\boxed{\cos\angle F_1PF_2=\frac35},\qquad
\angle F_1PF_2\ne90^\circ.
\]
而 $k_E=-1/2$、$k_H=2$，故切线交角是
\[
\boxed{\angle(t_E,t_H)=90^\circ}.
\]
\begin{center}
% P=(0,0) is a translated copy of (2,3); focal rays (-4,-3),(0,-3).
\begin{tikzpicture}[x=.95cm,y=.95cm,>=stealth,font=\small]
\draw[gray,thick] (0,0)--(-4,-3) (0,0)--(0,-3);
\draw[gray,dotted] (0,0)--(0,1.6);
\draw[orange!80!black,thick,dashed] (-1.7,-3.4)--(.9,1.8);
\draw[teal!85!black,thick,dashed] (-3,1.5)--(2,-1);
\draw[orange!80!black] ({.7*cos(216.86989765)},{.7*sin(216.86989765)})
 arc[start angle=216.86989765,end angle=243.43494882,radius=.7];
\draw[orange!80!black] ({.85*cos(243.43494882)},{.85*sin(243.43494882)})
 arc[start angle=243.43494882,end angle=270,radius=.85];
\draw[teal!85!black] (0,1.05) arc[start angle=90,end angle=153.43494882,radius=1.05];
\draw[teal!85!black] ({1.2*cos(153.43494882)},{1.2*sin(153.43494882)})
 arc[start angle=153.43494882,end angle=216.86989765,radius=1.2];
\draw (.268,-.134)--(.402,.134)--(.134,.268);
\fill (0,0) circle (2pt) node[above right=7pt] {$P$};
\fill (-4,-3) circle (2pt) node[below=5pt] {$F_1$};
\fill (0,-3) circle (2pt) node[below right=5pt] {$F_2$};
\node[orange!80!black,above=5pt] at (.9,1.8) {$t_H$：内角平分线};
\node[teal!85!black,anchor=east] at (-2.5,1.7) {$t_E$：外角平分线};
\node at (-1,-4.1) {内、外角平分线夹角为 $90^\circ$};
\end{tikzpicture}


{\small 例1的焦角满足 $\cos\theta=3/5$；直角标记属于两条切线。}
\end{center}
\textbf{方法迁移：}出现“夹角”先问：到底是\emph{焦角}还是\emph{切线交角}？
前者依赖离心率，后者在共焦条件下恒定为直角。
\end{examplebox}


\clearpage
\begin{examplebox}
\textbf{例4（两步判断共焦：把平方关系逆向使用）。}

已知两条同中心、同坐标轴的曲线
\[
 E:\ \frac{x^2}{36}+\frac{y^2}{20}=1,\qquad
 H:\ \frac{x^2}{9}-\frac{y^2}{7}=1.
\]
不求交点，判断它们在公共点处是否正交。

\textbf{第一步（辨认参数）：}
$a^2=36,b^2=20,A^2=9,B^2=7$，所以
\[
 a^2-A^2=36-9=27,\qquad b^2+B^2=20+7=27.
\]
\textbf{第二步（判断焦点）：}
\[
 c_E^2=36-20=16,\qquad c_H^2=9+7=16.
\]
两条曲线共享焦点 $(\pm4,0)$，因此所有公共点处切线都垂直。
\textbf{迁移：}标准式已共轴时，比较“平方差与平方和”
能够更快看出是否共焦；\emph{无须先求四个交点}。
\end{examplebox}

\clearpage
\begin{examplebox}
\textbf{高价值变式（去掉“共焦”这个关键条件）。}

椭圆 $E:x^2/16+y^2/12=1$ 与双曲线
$H_1:x^2/2-y^2/9=1$ 都经过 $P=(2,3)$。
这两条曲线在 $P$ 处仍然正交吗？请说明理由。

\textbf{先独立判断（建议限时2分钟）：}不要把“同轴”误看作“共焦”。

\textbf{解：}椭圆 $c_E^2=4$，双曲线 $c_{H_1}^2=2+9=11$，
所以\textbf{不同焦}。切线斜率
\[
k_E=-\frac12,\qquad
k_{H_1}=\frac{9x}{2y}\bigg|_{(2,3)}=3.
\]
于是
\[
\boxed{k_Ek_{H_1}=-\frac32\ne-1},
\]
故\textbf{两切线不垂直}。
\begin{center}
% Same coordinate scales/ranges. H0:x²-y²/3=1; H1:x²/2-y²/9=1.
% Both pass through P=(2,3), slopes 2 and 3 respectively.
\begin{tikzpicture}[x=.65cm,y=.65cm,font=\small]
\foreach \offset/\case/\title in {0/0/共焦：正交,10/1/非共焦：不正交} {
\begin{scope}[shift={(\offset,0)}]
\node at (1.3,5.8) {\title};
\draw[blue!75!black,thick,domain=10:105,samples=80,variable=\t]
 plot ({4*cos(\t)},{sqrt(12)*sin(\t)});
\ifnum\case=0
\draw[red!75!black,thick,domain=.2:5,samples=80,variable=\t]
 plot ({sqrt(1+\t*\t/3)},\t);
\draw[orange!80!black,dashed,thick] (1,1)--(3,5);
\draw (2.358,2.821)--(2.537,3.179)--(2.179,3.358);
\node at (.5,5.1) {$90^\circ$};
\node at (1.4,-.65) {$k_Ek_H=-1$};
\else
\draw[red!75!black,thick,domain=.2:5,samples=80,variable=\t]
 plot ({sqrt(2+2*\t*\t/9)},\t);
\draw[orange!80!black,dashed,thick] (1.3,.9)--(2.7,5.1);
\draw[gray] ({2+.7*cos(71.56505118)},{3+.7*sin(71.56505118)})
 arc[start angle=71.56505118,end angle=153.43494882,radius=.7];
\node at (.5,5.1) {$\approx81.9^\circ$};
\node at (1.4,-.65) {$k_Ek_{H_1}=-\tfrac32$};
\fi
\draw[teal!85!black,dashed,thick] (-1,4.5)--(4.5,1.75);
\fill (2,3) circle (2.5pt);
\node[anchor=east] at (1.1,2.35) {$P(2,3)$};
\end{scope}}
\end{tikzpicture}


{\small 同尺度对照：同过 $P$，改变双曲线参数后，切线不再垂直。}
\end{center}
\textbf{反例价值：}两条曲线有相同中心、焦轴，甚至有公共点，
仍不能替代“两个焦点完全相同”这一条件。
\end{examplebox}
